% addendum-mechanical.tex
% Document 2 - Technical Discovery Addendum: Mechanical
%
% Build (from this folder, requires XeLaTeX -- see README.md):
%   xelatex -interaction=nonstopmode addendum-mechanical.tex
%   xelatex -interaction=nonstopmode addendum-mechanical.tex   (rerun for TOC/page refs)
%
% Issued alongside the Pre-Project Intake and Discovery Form (Document 1)
% only when mechanical engineering assistance was indicated. See
% docs/Engineering_Project_Intake_Plan.md for the full workflow.

\documentclass[11pt]{article}

\usepackage{../../shared/intake-style}
\usepackage{../../shared/form-components}

\SetDocIdentity{Technical Discovery Addendum -- Mechanical}{ENG-FORM-02M}{v0.1-draft}
\SetFieldPrefix{addmech}

% Replace these placeholders once branding and markings are confirmed
% (see "Remaining Organization-Specific Decisions" in the intake plan).
\SetCompanyName{[Engineering Center Name]}
\SetMarkingText{FOR OFFICIAL USE ONLY}

\SetDocPurpose{%
This addendum collects the mechanical-engineering detail that the
Pre-Project Intake and Discovery Form (Document 1) does not ask for. Fill
it out only if mechanical assistance was indicated on that form. Its
purpose is to give the center's mechanical engineers enough detail about
existing equipment, loads, materials, and constraints to scope the work
accurately. It is fine to mark an item Unknown -- the center will help
fill in technical gaps during discussion.%
}

\begin{document}

\begin{Form}

\CoverPage

\DocumentControlBlock

\LinkedIntakeBlock

\IntakeSection{Equipment and Systems Involved}
\TextArea{EquipDescribe}{Equipment / system description (type, make, model, age if known)}{\linewidth}{2cm}\\[8pt]
\TextArea{EquipDrawings}{Existing mechanical drawings, models, or manuals available}{\linewidth}{1.5cm}

\IntakeSection{Loads, Motion, and Thermal Considerations}
\TextArea{LoadsKnown}{Known loads, forces, speeds, or motion requirements}{\linewidth}{1.5cm}\\[8pt]
\TextArea{ThermalKnown}{Known thermal, pressure, or fluid-flow considerations}{\linewidth}{1.5cm}

\IntakeSection{Materials and Environment}
\TextArea{MaterialsKnown}{Known or required materials of construction}{\linewidth}{1.2cm}\\[8pt]
\TextArea{EnvExposure}{Operating environment (temperature, corrosion, vibration, outdoor exposure, etc.)}{\linewidth}{1.5cm}

\IntakeSection{Applicable Standards}
\Guidance{List any mechanical codes or standards already known to apply.
Leave Confirmed By / Status blank if not yet reviewed by a discipline
lead -- this records applicability, not compliance.}\\[4pt]
\StandardsTable{MechStd}

\IntakeSection{Site and Access Constraints}
\TextArea{SiteAccess}{Site access, outage windows, or safety constraints specific to this equipment}{\linewidth}{1.5cm}

\IntakeSection{Assumptions and Open Questions}
\OpenItemsTable{MechOpen}

\AcknowledgementBlock{This addendum supplements the Pre-Project Intake and
Discovery Form and reflects current understanding only. It does not commit
\CompanyName's resources and is not authorization to begin engineering
work or to charge project time.}

\end{Form}

\end{document}
